\begin{abstract}
\noindent
Selective visual attention requires an animal to do two different things at
once: decide, moment by moment, \emph{what to do} on the basis of incoming
sensory evidence, and estimate \emph{how valuable} the current situation is, the
quantity that supplies the learning signal for the first computation. In the
primate brain these functions are carried by anatomically distinct but tightly
interacting circuits --- a sensory-grounded priority and commitment system in
parietal, frontal and midbrain nodes (LIP, FEF, superior colliculus) and a
value/prediction-error system in the cortico-basal-ganglia loop. Here we ask
whether the \emph{organisation} of these two computations, not merely their
presence, is a determinant of whether attention-like behaviour can be learned at
all. We study a recurrent vision-transformer agent trained purely from sparse
reward to perform a spatially cued orientation-change-detection task modelled on
primate covert-attention paradigms. We split the agent into two recurrent
cross-attention pathways --- a \emph{$\mu$ pathway} that reads the live image and
drives the policy (when to declare a change) and a \emph{$q$ pathway} that reads
accumulated memory and drives the value estimate --- and we vary a single
structural property: whether the pathways \emph{cross-talk} through a shared
recurrent memory or run as isolated modules. The result is categorical. With
cross-talk the agent becomes a calibrated psychophysical detector that reproduces
the hallmark reliability-scaled cueing signature of primate spatial attention
(\TODO{acc / hit / FA / RT placeholders}). Severing the coupling --- giving each
pathway a private memory --- abolishes learning entirely: the agent never
declares a change (\TODO{hit$=$0}). A second control that retains the coupling but
routes the memory pathway to the policy collapses identically, establishing two
jointly necessary conditions: the pathways must share a substrate, and the
image-grounded pathway must command the decision. We trace the failure to credit
assignment --- a critic that evaluates a different state than the actor occupies
cannot teach it --- the same principle that binds dopaminergic value signals to
cortico-striatal action learning. We further show the architecture generalises to
larger stimulus arrays (a nine-stimulus set-size manipulation; \TODO{set-size
$\times$ validity result}) and to environments with competing distractors
(\TODO{distractor result}), and we use the trained model to generate falsifiable
predictions --- a representational and causal dissociation between the policy and
value pathways --- about what should be decodable from, and disrupted by
perturbing, the corresponding neural circuits. The work illustrates how a
reward-trained neural network, constrained to be interpretable and manipulable,
can serve as a generative source of mechanistic hypotheses for systems
neuroscience.
\end{abstract}
